\documentclass[10pt,a4paper]{amsart}
\usepackage{amsmath,amssymb,amsthm,graphicx}
\DeclareMathOperator{\ev}{ev}
\DeclareMathOperator{\gr}{gr}
\DeclareMathOperator{\spn}{span}
% Minimal class-compatibility preamble for a reader excerpt of
%   Bhimarthi Ravinder, "Demazure Modules, Chari-Venkatesh Modules and Fusion Products",
%   SIGMA 10 (2014), 110, DOI 10.3842/SIGMA.2014.110.
%
% The author's pinned TeX file declares no theorem environments of its own: they are
% supplied by the SIGMA publisher class sigma.cls, which is not installed here.  The
% block below is that class's theorem-environment block (same environment names and
% styles), copied mechanically so that the retained author text compiles unchanged.
% No mathematical macro or statement is introduced by this file.
\newtheorem{theorem}{Theorem}
\newtheorem*{theorem*}{Theorem}
\newtheorem{lemma}{Lemma}
\newtheorem*{lemma*}{Lemma}
\newtheorem{corollary}{Corollary}
\newtheorem*{corollary*}{Corollary}
\newtheorem{proposition}{Proposition}
\newtheorem*{proposition*}{Proposition}
\newtheorem{conjecture}{Conjecture}
\newtheorem*{conjecture*}{Conjecture}
{\theoremstyle{definition} \newtheorem{definition}{Definition}
\newtheorem*{definition*}{Definition}
\newtheorem{example}{Example}
\newtheorem*{example*}{Example}
\newtheorem{remark}{Remark}
\newtheorem*{remark*}{Remark}
\newtheorem{note}{Note}
\newtheorem*{note*}{Note}
}
\usepackage{cleveref}
\crefname{originalreference}{original reference}{original references}
\Crefname{originalreference}{Original reference}{Original references}
\makeatletter
\newlabel{f1}{{f1}{source}{}{sourcepointer-1}{}}
\newlabel{f1@cref}{{[originalreference][0][]f1}{[1][0][]source}{}{}{}}
\makeatother
\title{Annihilation in a fusion product from annihilation in the factors}
\author{Bhimarthi Ravinder}
\date{doi:10.3842/SIGMA.2014.110; CC BY 4.0}
\begin{document}
\maketitle
\noindent\textit{Editorial note.} Preamble provenance: the article keeps its theorem environments in the SIGMA publisher class and declares three operators of its own in the preamble of the pinned tex. This card therefore uses two verbatim preamble spans: the operator declarations from the pinned source (lines 9-11, required because the counted proof writes the associated graded with the operator declared there) and the staged retained-preamble.tex theorem-environment block (lines 1-28). Nothing was added, renamed or redefined, and extra\_packages stays empty; the wrapper class and base packages are supplied by the assembler.\par\medskip
\noindent\textit{Editorial note.} All delivered mathematical text is the author's own, copied from the pinned article source sources/sigma-2014-110-demazure-fusion/pinned-original-source.tex; no statement or proof step was rewritten, reordered or removed.\par\medskip
\noindent\textit{Editorial note.} Preamble provenance: the author's tex declares no theorem environments of its own (they come from the SIGMA publisher class); the staged sources/sigma-2014-110-demazure-fusion/retained-preamble.tex reproduces that class's theorem-environment block verbatim and is used as the preamble, with the wrapper class and base packages supplied by the assembler.\par\medskip
\noindent\textit{Editorial note.} Bibliography provenance: neither the counted statement nor its proof contains a citation (this lane rescans the whole span before writing the card), so the card carries no bibliography and no entry was reconstructed.\par\medskip
\noindent\textit{Editorial note.} Selection provenance: this card was opened under the file-based control contract (CONTROL.md), with the mechanical pre-card gate run on its statement and proof interval before the card was written.\par\medskip
\noindent\textit{Allowed standard prerequisites (editorial).}\begin{itemize}
\item The fusion product of cyclic graded modules, defined through the associated graded of the tensor product with respect to the filtration by the powers of a parameter. The author's section 4.1, in which the fusion product is set up before the counted statement. Fixes the product in which the annihilation is asserted.
\item The identification of the operator acting on the tensor product of the factors with the product of the shifted operators acting on the generators. Invoked inside the counted proof at source line 568 as a pointer to the lemma labelled f1 (a navigation-pointer stub). Computes the action of the twisted operator on the class of the tensor of the generators.
\end{itemize}\medskip
\ifdefined\hypertarget\hypertarget{sourcepointer-1}{}\fi
\noindent\textit{Source reference.} Original reference f1: sources/sigma-2014-110-demazure-fusion/pinned-original-source.tex, source lines 568--568. Navigation pointer only: the counted text refers to the item f1 of this same article, which is not part of this card. The pointer carries the item's own author label as its number; its content is not reproduced here.\par\medskip
\begin{lemma}
\label{f2}
Given $1\leq i\leq m$, let $V_i$ be a~finite-dimensional cyclic graded $\mathfrak{g}[t]$-module generated by $v_i$, and
$s_i\in\mathbb{Z}_{\geq0}$.
Let $x\in\mathfrak{g}$.
If $(x\otimes t^{s_i})v_i=0$, $\forall\, 1\leq i\leq m$ then $(x\otimes t^{s_1+\dots+s_m}) v_1*\dots*v_m=0$.
\end{lemma}

\begin{proof}
Let $z_1,\dots,z_m$ be distinct complex numbers and let $\mathbf{V}$ as above.
By using Lemma~\ref{f1}, we get the following equality in $\gr \mathbf{V}$,
\begin{gather*}
\big(x\otimes t^{s_1+\cdots+s_m}\big) \big((v_1\otimes\cdots\otimes v_m)+F^{-1}\mathbf{V}\big)
\\
\qquad{}
=\big(\big(x\otimes (t-z_1)^{s_1}\cdots(t-z_m)^{s_m}\big)v_1\otimes\dots\otimes v_m\big) + F^{s_1+\cdots+s_m-1}\mathbf{V}.
\end{gather*}
Now the proof follows by the def\/inition of the fusion product.
\end{proof}

\end{document}
